\documentclass[11pt, a4paper]{article}

% --- UNIVERSAL PREAMBLE BLOCK ---
\usepackage[a4paper, top=2.5cm, bottom=2.5cm, left=2cm, right=2cm]{geometry}
\usepackage{fontspec}
\usepackage[english, bidi=basic, provide=*]{babel}

% Set default/Latin font to Sans Serif in the main (rm) slot
\babelfont{rm}{Noto Sans}

\usepackage{amsmath}
\usepackage{amssymb}
\usepackage{enumitem}
\usepackage{xlop} % For arithmetic vertical alignment
\usepackage{fancyhdr}
\usepackage{xcolor}

\newcommand{\answerline}[1]{\textcolor{red}{\textbf{#1}}}
\newcommand{\working}[1]{\vspace{10pt}\textcolor{blue}{\small \textbf{Working/Steps:}\\ #1}}

% Page numbering and footer
\pagestyle{fancy}
\fancyhf{}
\renewcommand{\headrulewidth}{0pt}
\fancyfoot[C]{\thepage}
\fancyfoot[L]{\small Singapore Math Level 4A \& 4B (SOLUTION)}

\begin{document}

\begin{center}
    \framebox{\parbox{0.9\textwidth}{\centering \huge \textbf{Unit 3: WHOLE NUMBERS (PART 3) - SOLUTIONS}}}
\end{center}

\vspace{1cm}

\textbf{Solve the following problems. Show your work.}

\begin{enumerate}[label=\arabic*.]
    \item \begin{tabular}[t]{c r}
          & 4 1 2 \\
        $\times$ & 4 \\ \hline
          & \answerline{1, 6 4 8} \\ \hline
    \end{tabular}
    \hfill
    \setcounter{enumi}{5}
    \item \begin{tabular}[t]{c r}
          & 5, 3 1 7 \\
        $\times$ & 6 \\ \hline
          & \answerline{3 1, 9 0 2} \\ \hline
    \end{tabular}

    \vspace{1cm}
    \setcounter{enumi}{1}
    \item \begin{tabular}[t]{c r}
          & 5 4 7 \\
        $\times$ & 2 \\ \hline
          & \answerline{1, 0 9 4} \\ \hline
    \end{tabular}
    \hfill
    \setcounter{enumi}{6}
    \item \begin{tabular}[t]{c r}
          & 2, 0 1 1 \\
        $\times$ & 8 \\ \hline
          & \answerline{1 6, 0 8 8} \\ \hline
    \end{tabular}

    \vspace{1cm}
    \setcounter{enumi}{2}
    \item \begin{tabular}[t]{c r}
          & 6 1 0 \\
        $\times$ & 5 \\ \hline
          & \answerline{3, 0 5 0} \\ \hline
    \end{tabular}
    \hfill
    \setcounter{enumi}{7}
    \item \begin{tabular}[t]{c r}
          & 6, 0 2 8 \\
        $\times$ & 9 \\ \hline
          & \answerline{5 4, 2 5 2} \\ \hline
    \end{tabular}

    \vspace{1cm}
    \setcounter{enumi}{3}
    \item \begin{tabular}[t]{c r}
          & 9 3 5 \\
        $\times$ & 3 \\ \hline
          & \answerline{2, 8 0 5} \\ \hline
    \end{tabular}
    \hfill
    \setcounter{enumi}{8}
    \item \begin{tabular}[t]{c r}
          & 1, 5 2 6 \\
        $\times$ & 5 \\ \hline
          & \answerline{7, 6 3 0} \\ \hline
    \end{tabular}

    \vspace{1cm}
    \setcounter{enumi}{4}
    \item \begin{tabular}[t]{c r}
          & 1 0 9 \\
        $\times$ & 7 \\ \hline
          & \answerline{7 6 3} \\ \hline
    \end{tabular}
    \hfill
    \setcounter{enumi}{9}
    \item \begin{tabular}[t]{c r}
          & 8, 4 3 7 \\
        $\times$ & 8 \\ \hline
          & \answerline{6 7, 4 9 6} \\ \hline
    \end{tabular}
\end{enumerate}

\pagebreak

\begin{enumerate}[label=\arabic*.]
    \setcounter{enumi}{10}
    \item \begin{tabular}[t]{c r}
          & 4 6 \\
        $\times$ & 1 8 \\ \hline
          & \answerline{8 2 8} \\ \hline
    \end{tabular}
    \hfill
    \setcounter{enumi}{15}
    \item \begin{tabular}[t]{c r}
          & 1 2 6 \\
        $\times$ & 5 0 \\ \hline
          & \answerline{6, 3 0 0} \\ \hline
    \end{tabular}

    \vspace{1.5cm}
    \setcounter{enumi}{11}
    \item \begin{tabular}[t]{c r}
          & 3 5 \\
        $\times$ & 2 0 \\ \hline
          & \answerline{7 0 0} \\ \hline
    \end{tabular}
    \hfill
    \setcounter{enumi}{16}
    \item \begin{tabular}[t]{c r}
          & 6 2 5 \\
        $\times$ & 7 3 \\ \hline
          & \answerline{4 5, 6 2 5} \\ \hline
    \end{tabular}

    \vspace{1.5cm}
    \setcounter{enumi}{12}
    \item \begin{tabular}[t]{c r}
          & 6 7 \\
        $\times$ & 3 6 \\ \hline
          & \answerline{2, 4 1 2} \\ \hline
    \end{tabular}
    \hfill
    \setcounter{enumi}{17}
    \item \begin{tabular}[t]{c r}
          & 6 1 9 \\
        $\times$ & 2 4 \\ \hline
          & \answerline{1 4, 8 5 6} \\ \hline
    \end{tabular}

    \vspace{1.5cm}
    \setcounter{enumi}{13}
    \item \begin{tabular}[t]{c r}
          & 9 1 \\
        $\times$ & 2 7 \\ \hline
          & \answerline{2, 4 5 7} \\ \hline
    \end{tabular}
    \hfill
    \setcounter{enumi}{18}
    \item \begin{tabular}[t]{c r}
          & 2 8 1 \\
        $\times$ & 5 3 \\ \hline
          & \answerline{1 4, 8 9 3} \\ \hline
    \end{tabular}

    \vspace{1.5cm}
    \setcounter{enumi}{14}
    \item \begin{tabular}[t]{c r}
          & 8 9 \\
        $\times$ & 1 6 \\ \hline
          & \answerline{1, 4 2 4} \\ \hline
    \end{tabular}
    \hfill
    \setcounter{enumi}{19}
    \item \begin{tabular}[t]{c r}
          & 3 8 0 \\
        $\times$ & 3 6 \\ \hline
          & \answerline{1 3, 6 8 0} \\ \hline
    \end{tabular}
\end{enumerate}

\pagebreak

\textbf{Solve the following problems. Show your work.}

\begin{enumerate}[label=\arabic*.]
    \setcounter{enumi}{20}
    \item $5 \overline{) 1, 3 5 5}$ \quad \answerline{Answer: 271}
    \hfill
    \setcounter{enumi}{23}
    \item $8 \overline{) 6, 0 8 8}$ \quad \answerline{Answer: 761}

    \vspace{2cm}
    \setcounter{enumi}{21}
    \item $3 \overline{) 4, 8 2 7}$ \quad \answerline{Answer: 1,609}
    \hfill
    \setcounter{enumi}{24}
    \item $6 \overline{) 1, 4 5 8}$ \quad \answerline{Answer: 243}

    \vspace{2cm}
    \setcounter{enumi}{22}
    \item $2 \overline{) 9, 8 0 4}$ \quad \answerline{Answer: 4,902}
    \hfill
    \setcounter{enumi}{25}
    \item $3 \overline{) 1, 1 3 1}$ \quad \answerline{Answer: 377}
\end{enumerate}

\vspace{1cm}
\textbf{Fill in each blank with the correct answer.}

\begin{enumerate}[label=\arabic*.]
    \setcounter{enumi}{26}
    \item $18 \times 20 = \text{\answerline{18}} \times \text{\answerline{2}} \text{ tens} = \text{\answerline{36}} \text{ tens} = \text{\answerline{360}}$
    \item $69 \times 40 = \text{\answerline{69}} \times \text{\answerline{4}} \text{ tens} = \text{\answerline{276}} \text{ tens} = \text{\answerline{2,760}}$
    \item $98 \times 30 = \text{\answerline{98}} \times \text{\answerline{3}} \text{ tens} = \text{\answerline{294}} \text{ tens} = \text{\answerline{2,940}}$
    \item $53 \times 60 = \text{\answerline{53}} \times \text{\answerline{6}} \times 10 = \text{\answerline{318}} \times 10 = \text{\answerline{3,180}}$
    \item $77 \times 90 = \text{\answerline{77}} \times \text{\answerline{9}} \times 10 = \text{\answerline{693}} \times 10 = \text{\answerline{6,930}}$
    \item $42 \times 80 = \text{\answerline{42}} \times \text{\answerline{8}} \times 10 = \text{\answerline{336}} \times 10 = \text{\answerline{3,360}}$
\end{enumerate}

\pagebreak

\begin{enumerate}[label=\arabic*.]
    \setcounter{enumi}{32}
    \item $4,569 \div 8 = \text{\answerline{571}} \text{ R } \text{\answerline{1}}$
    \item $1,348 \div 5 = \text{\answerline{269}} \text{ R } \text{\answerline{3}}$
    \item $4,240 \div 7 = \text{\answerline{605}} \text{ R } \text{\answerline{5}}$
    \item $3,134 \div 4 = \text{\answerline{783}} \text{ R } \text{\answerline{2}}$
    \item $9,381 \div 9 = \text{\answerline{1,042}} \text{ R } \text{\answerline{3}}$

    \vspace{1cm}
    \item $59 \times 17 \approx \text{\answerline{60}} \times \text{\answerline{20}} = \text{\answerline{1,200}}$
    \item $623 \times 55 \approx \text{\answerline{600}} \times \text{\answerline{60}} = \text{\answerline{36,000}}$
    \item $917 \times 31 \approx \text{\answerline{900}} \times \text{\answerline{30}} = \text{\answerline{27,000}}$
    \item $237 \div 5 \approx \text{\answerline{250}} \div 5 = \text{\answerline{50}}$
    \item $4,927 \div 7 \approx \text{\answerline{4,900}} \div 7 = \text{\answerline{700}}$
    \item $8,080 \div 9 \approx \text{\answerline{8,100}} \div 9 = \text{\answerline{900}}$
\end{enumerate}

\pagebreak

\begin{enumerate}[label=\arabic*.]
    \setcounter{enumi}{43}
    \item (a) Multiply 84 by 52. \hfill \text{\answerline{4,368}} \\
          (b) Estimate the value of $84 \times 52$. \hfill \text{\answerline{80 \times 50 = 4,000}} \\
          (c) State if your actual answer is reasonable. \hfill \text{\answerline{Yes, 4,368 is close to 4,000}}

    \item (a) Multiply 409 by 93. \hfill \text{\answerline{38,037}} \\
          (b) Estimate the value of $409 \times 93$. \hfill \text{\answerline{400 \times 90 = 36,000}} \\
          (c) State if your actual answer is reasonable. \hfill \text{\answerline{Yes, 38,037 is close to 36,000}}

    \item (a) Multiply 976 by 38. \hfill \text{\answerline{37,088}} \\
          (b) Estimate the value of $976 \times 38$. \hfill \text{\answerline{1,000 \times 40 = 40,000}} \\
          (c) State if your actual answer is reasonable. \hfill \text{\answerline{Yes, 37,088 is close to 40,000}}

    \item (a) Divide 43 by 8. \hfill \text{\answerline{5 R 3}} \\
          (b) Estimate the value of $43 \div 8$. \hfill \text{\answerline{40 \div 8 = 5}} \\
          (c) State if your actual answer is reasonable. \hfill \text{\answerline{Yes}}

    \item (a) Divide 538 by 6. \hfill \text{\answerline{89 R 4}} \\
          (b) Estimate the value of $538 \div 6$. \hfill \text{\answerline{540 \div 6 = 90}} \\
          (c) State if your actual answer is reasonable. \hfill \text{\answerline{Yes}}

    \item (a) Divide 8,765 by 4. \hfill \text{\answerline{2,191 R 1}} \\
          (b) Estimate the value of $8,765 \div 4$. \hfill \text{\answerline{8,800 \div 4 = 2,200}} \\
          (c) State if your actual answer is reasonable. \hfill \text{\answerline{Yes}}
\end{enumerate}

\pagebreak

\textbf{Solve the following story problems. Show your work in the space below.}

\begin{enumerate}[label=\arabic*.]
    \setcounter{enumi}{49}
    \item There are 255 balloons in a package. How many balloons are there in a dozen packages?
    \working{1 dozen = 12 \\ 255 $\times$ 12 = 3,060 \\ \answerline{There are 3,060 balloons.}}

    \vspace{2cm}

    \item A factory makes 2,275 watches in a week. How many watches does it make in 3 days?
    \working{1 week = 7 days \\ Watches made in 1 day: 2,275 $\div$ 7 = 325 \\ Watches made in 3 days: 325 $\times$ 3 = 975 \\ \answerline{The factory makes 975 watches in 3 days.}}

    \vspace{2cm}

    \item A shirt costs \$253 and a tie costs \$78. David bought 4 shirts and 3 ties. How much did he spend altogether?
    \working{Cost of 4 shirts = \$253 $\times$ 4 = \$1,012 \\ Cost of 3 ties = \$78 $\times$ 3 = \$234 \\ Total cost = \$1,012 + \$234 = \$1,246 \\ \answerline{He spent \$1,246 altogether.}}

    \vspace{2cm}

    \item A purse costs 4 times as much as a dress. If the purse costs \$276, how much does Anna spend on the purse and 3 dresses?
    \working{Cost of 1 dress = \$276 $\div$ 4 = \$69 \\ Cost of 3 dresses = \$69 $\times$ 3 = \$207 \\ Total spent = Purse + 3 Dresses = \$276 + \$207 = \$483 \\ \answerline{Anna spends \$483.}}

\end{enumerate}

\pagebreak

\begin{enumerate}[label=\arabic*.]
    \setcounter{enumi}{53}
    \item Jason had some marbles. He gave 35 marbles to each of his 4 brothers and still had 219 marbles left. How many marbles did Jason have to begin with?
    \working{Marbles given to 4 brothers = 35 $\times$ 4 = 140 \\ Total marbles he had carefully given away + left over: \\ 140 + 219 = 359 \\ \answerline{He had 359 marbles to begin with.}}

    \vspace{4cm}

    \item 114 men and 686 women went to a concert. Each ticket cost \$17. How much money was collected in all?
    \working{Total number of people = 114 + 686 = 800 \\ Total money collected = 800 $\times$ \$17 = \$13,600 \\ \answerline{\$13,600 was collected in all.}}

    \vspace{4cm}

    \item A baker bakes 840 loaves of bread in a day. How many loaves of bread will he bake in 6 weeks?
    \working{6 weeks = 6 $\times$ 7 days = 42 days \\ Total loaves = 840 $\times$ 42 = 35,280 \\ \answerline{He will bake 35,280 loaves of bread.}}

\end{enumerate}

\pagebreak

\begin{enumerate}[label=\arabic*.]
    \setcounter{enumi}{56}
    \item Cecilia has 896 stickers. She gives 50 stickers to seven of her friends. She sorts the remaining stickers equally into three albums. How many stickers are there in each album?
    \working{Stickers given away: 50 $\times$ 7 = 350 \\ Remaining stickers: 896 - 350 = 546 \\ Stickers per album: 546 $\div$ 3 = 182 \\ \answerline{There are 182 stickers in each album.}}

    \vspace{4cm}

    \item There were 400 pieces of candy in a package. The principal of a school bought 25 packages for 2,000 children on Halloween.
    
    \begin{enumerate}[label=(\alph*), leftmargin=*]
        \item How many pieces of candy did the principal buy altogether?
        \working{400 $\times$ 25 = 10,000 \\ \answerline{The principal bought 10,000 pieces.}}
        
        \item If each child was given 7 pieces, how many more packages were needed?
        \working{Candy needed for all children = 2,000 $\times$ 7 = 14,000 \\ Shortage of candy = 14,000 - 10,000 = 4,000 \\ Additional packages needed = 4,000 $\div$ 400 = 10 \\ \answerline{10 more packages were needed.}}
    \end{enumerate}

\end{enumerate}

\pagebreak

\begin{enumerate}[label=\arabic*.]
    \setcounter{enumi}{58}
    \item Kimi is 16 years old and her mother is 44 years old this year. How many years ago was Kimi's mother five times as old as Kimi?
    \working{Difference in age is always constant: 44 - 16 = 28 years. \\ In the past, if Mother was 5 times older, then: \\ Mother = 5 units, Kimi = 1 unit. \\ Difference = 4 units. \\ Therefore, 4 units = 28 $\Rightarrow$ 1 unit = 7. \\ Kimi was 7 years old at that time. \\ Currently, Kimi is 16. \\ Number of years ago: 16 - 7 = 9. \\ \answerline{It was 9 years ago.}}

    \vspace{4cm}

    \item A stereo costs \$328. A television set costs four times as much as the stereo. Mr. Simon buys the stereo and the television set and pays for them in eight monthly installments. How much must he pay for them each month?
    \working{Cost of television = \$328 $\times$ 4 = \$1,312 \\ Total cost = Stereo + Television = \$328 + \$1,312 = \$1,640 \\ Monthly installment = \$1,640 $\div$ 8 = \$205 \\ \answerline{He must pay \$205 each month.}}

    \vspace{4cm}

    \item Michael had \$3,600. After spending \$320, he still had twice as much as Cynthia. Find the total amount of money they had in the beginning.
    \working{Money Michael had left = \$3,600 - \$320 = \$3,280 \\ Michael's remainder is twice Cynthia's $\Rightarrow$ Cynthia had \$3,280 $\div$ 2 = \$1,640 \\ Total amount in the beginning = Michael's \$3,600 + Cynthia's \$1,640 = \$5,240 \\ \answerline{The total amount was \$5,240.}}

\end{enumerate}

\pagebreak

\begin{enumerate}[label=\arabic*.]
    \setcounter{enumi}{61}
    \item Luis bought a book and four identical pens for \$12. Carlos bought the same book and two similar pens. Carlos paid \$4 less than Luis. What was the cost of the book?
    \working{Luis paid for: 1 Book + 4 Pens = \$12 \\ Carlos paid for: 1 Book + 2 Pens = \$12 - \$4 = \$8 \\ The difference is 2 Pens, which costs \$4. \\ Cost of 1 Pen = \$4 $\div$ 2 = \$2. \\ Since Carlos paid \$8 for 1 Book + 2 Pens (which cost \$4), \\ the Book costs \$8 - \$4 = \$4. \\ \answerline{The cost of the book was \$4.}}

    \vspace{4cm}

    \item 250 adults and some children went to the zoo. The admission ticket for each adult was \$12 and the admission ticket for each child was \$9. If \$6,915 was collected for all the tickets, how many children went to the zoo?
    \working{Money collected from adults = 250 $\times$ \$12 = \$3,000 \\ Money collected from children = Total money - Adult money \\ = \$6,915 - \$3,000 = \$3,915 \\ Number of children = \$3,915 $\div$ \$9 = 435 \\ \answerline{There were 435 children.}}

    \vspace{4cm}

    \item Mr. Ortiz received a bonus. He gave \$2,000 to his wife and distributed an amount of money equally among his six children. He was left with \$1,350, which was \$400 more than the amount of money he gave to each child. How much was his bonus?
    \working{Money left = \$1,350. This is \$400 more than 1 child's share. \\ Amount given to 1 child = \$1,350 - \$400 = \$950 \\ Amount given to 6 children = \$950 $\times$ 6 = \$5,700 \\ Total bonus = Wife's share + Children's share + Amount left over \\ = \$2,000 + \$5,700 + \$1,350 = \$9,050 \\ \answerline{His bonus was \$9,050.}}

\end{enumerate}

\end{document}
